\documentclass{birkjour}
%%%%%%%%%%%%%%%%%%%%%%%%%%%%%%%%%%%%%%%
% Compile order:
% pdflatex -interaction=nonstopmode main_revised_v4
% bibtex main_revised_v4
% pdflatex -interaction=nonstopmode main_revised_v4
% pdflatex -interaction=nonstopmode main_revised_v4
%%%%%%%%%%%%%%%%%%%%%%%%%%%%%%%%%%%%%%%%%%%%%%%

\usepackage[T1]{fontenc}
\usepackage{tgtermes}
\usepackage[scaled=.92]{helvet}
\usepackage{amsthm,amsmath,amssymb}
\usepackage{multicol}
\usepackage[all,pdf]{xy}
\usepackage[utf8]{inputenc}
\usepackage{float}
\usepackage{mathrsfs}
\usepackage{tabularx}
\usepackage{multirow}
\usepackage[numbers]{natbib}
\usepackage[colorlinks,citecolor=blue,urlcolor=blue,linkcolor=blue]{hyperref}
\usepackage{subcaption}
\usepackage{stmaryrd}
\usepackage{comment}
\usepackage{proof}

\begin{document}

\title{Chin Yueh-Lin on Carroll's Paradox: Inference, Russell, and Logical Realism}
\author{Xinwen Liu}
\address{%
Institute of Philosophy, Chinese Academy of Social Sciences\\
Beijing 100732\\
PRC}
\email{liuxw-zxs@cass.org.cn}

\author{Rui Zhu}
\address{%
Level 3, Unit 16, Datang Hydropower Plant\\
West Triangle, Chengguan Town\\
Shiquan County, Shaanxi Province\\
PRC}
\email{raycaesar@126.com}

\subjclass{Primary 03A05; Secondary 01A25}
\keywords{theory of inference, Carroll's paradox, Chin Yueh-Lin, Russell, logical realism, foundational holism}
\date{}

\begin{abstract}
This paper reconstructs Chin Yueh-Lin's (1895--1984) discussion of Lewis Carroll's paradox in Chin's 1960 essay `On ``Therefore'''. The central claim is that Chin's response is Russellian in origin but anti-Russellian in philosophical development. Chin explicitly relates his own account of assertion to Russell's treatment of implication and inference in \textit{The Principles of Mathematics}, \S\,38. Yet he also insists that, unlike Russell, he does not strip assertion of its cognitive component. We argue that this difference is the key to the essay. Chin uses Carroll's paradox not to provide a solution superior to Russell's, nor to anticipate later proof-theoretic or inferentialist theories of logical normativity, but to defend the objectivity of inference while refusing to reduce inference to an abstract relation between unasserted propositions, a psychological transition, or a construction out of subjective materials. This reading also clarifies the relation between `On ``Therefore''', Chin's 1962 response to Zhou Liquan, \textit{The Philosophy of Russell}, and his later English critique of Russell's neutral monism. Across these texts, Chin's concern is continuous: inference, judgment, and scientific explanation are historically situated human activities, but they remain answerable to objective facts. We further compare this realist and historical account of inference with Gila Sher's foundational holism, cautiously and heuristically rather than genealogically. Finally, we argue that this reading challenges Hao Wang's ``two ideals'' interpretation of Chin: the vocabulary and systematic form of Chin's philosophy changed after 1949, but his concern with the objectivity of logic and inference persisted.
\end{abstract}
\maketitle

\section{Introduction}

Chin Yueh-Lin, or Jin Yuelin,\footnote{Throughout this paper, we use the Wade-Giles romanisation ``Chin Yueh-Lin'' for his pre-1949 English publications, e.g.~\cite{Chin1934}, to maintain historical and bibliographical consistency. However, following the widespread adoption of Pinyin after 1949, his later Chinese works in translation and post-1949 references are typically cited under ``Jin Yuelin.'' We retain ``Jin Yuelin'' in the bibliography for his post-1949 publications accordingly.} (1895--1984) was the most significant philosopher in modern China to have specialised in logic. His importance is not merely regional. During the 1930s and 1940s, he played a decisive role in establishing what historians of Chinese philosophy now recognise as the Tsinghua School of logic and philosophy.\footnote{During the 1930s and 1940s, the Department of Philosophy at Tsinghua University became the epicentre of logical and analytic philosophical studies in China. Under Chin's leadership, this academic community engaged deeply with the Western analytic tradition and mathematical logic. For a comprehensive history of the Tsinghua School, see Hu~\cite{Hu2005} and Vrhovski~\cite{vrhovski2021qinghua}.} Among his students at Tsinghua University was Hao Wang (1921--1995), later one of the leading mathematical logicians of the twentieth century. The intellectual route from Russell through Chin to Wang is therefore one of the major channels through which modern logic entered Chinese philosophy.

Chin's relation to Russell was never one of passive reception. He absorbed from early Russell and Moore a realist conception of facts, propositions, and logical relations; he also adopted Russellian logic as a model for philosophical clarity. Yet he gradually came to think that Russell's later epistemological programme endangered the very objectivity that Russell's early realism had seemed to secure. The most illuminating place to see this development is not, as one might expect, a purely technical logical work, but Chin's 1960 essay `On ``Therefore''', a Chinese text devoted in part to Lewis Carroll's 1895 dialogue ``What the Tortoise Said to Achilles.''

Among Chin's many publications in the philosophy of logic, `On ``Therefore''' is one of the most contested. It opens with the striking claim that ``therefore,'' and hence inference, is in some sense historically and socially situated. This claim was immediately controversial among Chin's students and later readers. But the opening slogan is not the deepest point of the essay. The deeper point is that Chin uses Carroll's paradox to defend the objectivity and indispensability of inference. The Tortoise appears to show that inference cannot get started; Russell's later constructional method appears to show that inference to objective entities can be replaced by logical constructions. Chin's answer is not that these are the same problem. It is that both threaten, in different ways, the role of inference as a bridge between thought and fact.

The textual basis for this claim is unusually explicit. In the eighth section of `On ``Therefore''', Chin directly refers to Russell's \textit{The Principles of Mathematics}, \S\,38. He notes that Russell had already attempted to strip assertion of its psychological component, and then marks his own difference from Russell: the present paper, Chin says, does not strip assertion of its cognitive component~\cite[p.~42]{Liu2018}. This passage shows that Chin's response to Carroll is not an independent reinvention of a solution already available in Russell. It is a self-conscious transformation of Russell's distinction between what is asserted and the act of asserting. The transformation is philosophical rather than merely technical. Chin accepts a Russellian distinction in order to block Carroll's regress, but then turns the question of assertion and inference against what he takes to be the individualistic and objectivist assumptions of Russell's later philosophy. By 'individualistic,' Chin means the treatment of inference as an act of an isolated knower, abstracted from historical and social conditions.

This paper therefore does not claim that Chin offered a solution to Carroll's paradox superior to Russell's. At the local level, his answer remains close to the Russellian distinction between implication and inference, or between the truth of a proposition and the act of asserting it. Nor do we claim that Chin anticipated later inferentialist or proof-theoretic accounts of logical normativity, that is, accounts which explain logical meaning or inferential entitlement in terms of rules of inference.\footnote{This later line of work includes, among others, Gentzen's proof theory, Prawitz's proof-theoretic semantics, Dummett's theory of the justification of deduction, Brandom's inferentialism, Schroeder-Heister's proof-theoretic semantics, and Boghossian's discussion of blind reasoning; see~\cite{Brandom1994,SchroederHeister2012,Boghossian2003}. Chin's project is not to provide a theory of this kind.} Those accounts ask how a reasoner is entitled to perform an inferential transition once the relevant rule or implication is recognised. Chin's concern is different: he wants to show that inference is not reducible to assertion, psychological transition, or subjective construction, but has an objective anchorage in facts.

Three claims organise the paper. First, Chin's response to Carroll is Russellian in origin but not merely Russellian in significance. He uses Russell's proposition/assertion distinction to resist Carroll's regress, while refusing Russell's attempt to remove the cognitive dimension of assertion. Second, Chin's discussion of Carroll prepares his later critique of Russell's constructional method. Russell's use of ``inference'' shifts between deductive and non-deductive contexts; Chin sees this instability and answers on both fronts, without claiming to have a single formal theory covering both. Third, Chin's appeal to ``historical facts'' and ``the possibility of thinking'' is best read as a realist and historical account of inference. This framework bears a structural affinity with Gila Sher's foundational holism, although Chin lacks Sher's formal model-theoretic apparatus. On this reading, Chin's later work does not simply break with his pre-1949 realism. It restates a persistent concern with the objectivity of logic in a changed philosophical vocabulary.

\section{Historical Background: `On ``Therefore''' and Its Self-Correction}
\label{sec:historical-background}

Chin Yueh-Lin's academic career divides roughly into two phases. Before the 1952 reorganisation of Chinese universities, he was based at Tsinghua University and, during the war years, at National Southwestern Associated University. Chin subsequently served as professor and chair of the Department of Philosophy at Peking University from 1952 to 1955. From 1955 until the end of his life, he served as the founder and first director of the Logic Section at the Institute of Philosophy, Chinese Academy of Social Sciences. `On ``Therefore''' appeared in \textit{Zhexue Yanjiu} (\textit{Philosophical Research}), the Institute's flagship journal, in 1960.

The paper opens with a sentence that has often overshadowed the rest of the essay: Chin writes that `` `therefore,' and hence inference, is basically of class characteristics''~\cite[p.~30]{Liu2017}. Taken literally, the sentence seems to threaten the universality of formal logic. This is how some readers understood it. Zhou Liquan (1921--2008), one of Chin's former students, argued that whether formal logic carried class characteristics was ``a question of life and death for logic''~\cite[p.~12]{Zhou2000}. Chin replied in 1962 in `On the Class Characteristics and Necessity of the Inference Form,' where he acknowledged serious defects in the 1960 paper:

\begin{quote}
The paper `On ``Therefore''' has many faults. Its scope is too large, and it raises so many issues that its main issue is much less prominent. The root of these faults lies in less clear thinking, whose cause is that old views have not been ruled out when it advocates new views. Its central idea is therefore somewhat immature.~\cite{Jin1962}\cite[vol.~4, p.~555]{Jin2014}
\end{quote}

This self-criticism is important, but it should not be read as a simple confession that the whole philosophical project of the essay was misguided. Two features deserve emphasis. First, Chin himself says that the main issue of the paper had become obscured. The controversial opening thesis is therefore not the whole argument. Second, his reference to ``old views'' not yet ruled out is revealing. In later memoir fragments written between 1981 and 1983 and published in 1995, Chin identifies those old views with his realism, and notes that his 1962 article on objective certainty and the laws of formal logic was among the three papers of which he was most satisfied with~\cite[vol.~4, pp.~897--899]{Jin2014}. What remained active in `On ``Therefore''' was not merely a residue to be eliminated, but the realist backbone of Chin's earlier philosophy.

The 1962 response to Zhou also clarifies what had remained unclear in the 1960 essay. Chin explicitly distinguishes logical implication from actual inference. In discussing syllogistic reasoning, he treats a first-figure syllogistic schema not as an inference itself, but as the formal ground of possible inferences: the schema expresses an implication relation, whereas actual inference occurs only when a thinker asserts the premises and moves to the conclusion~\cite{Jin1962}. Implication, on this account, has objective necessity; inference is a cognitive act performed by a situated thinker. Concrete inferences occur in determinate historical and social circumstances, whereas inference forms are objects of abstraction for formal logic. The point is not that formal validity is relative to class or circumstance. It is rather that actual inferential activity is not merely a relation among formal marks. It is a historically situated cognitive act that remains answerable to objective logical form and, in concrete cases, to the factual situation in which reasoning takes place.



This distinction stabilises the meaning of ``historical conditions'' in `On ``Therefore'''. This phrase has a layered meaning. It does not merely say that inference happens in time; it also marks the objective epistemic situation in which a cognitive transition is answerable to facts. In the weakest sense, it means simply that an inference is an event performed by a finite reasoner in a determinate situation. In a stronger epistemic sense, it means that the reasoner is not operating in a private mental vacuum but in an objective situation in which premises, facts, concepts, and practical interests are already structured. In the strongest realist sense, ``historical facts'' names the field of objective reality to which inference is answerable. These senses are related. Chin moves from the occurrence of inference to its epistemic situation, and from there to the realist claim that genuine inference must be situated in a world of objective facts. The phrase therefore should not be read as a causal-psychological explanation of why a person happens to infer, but as a realist description of the background within which an inferential act can count as more than a psychological transition.

This point also explains why `On ``Therefore''' should be read together with Chin's later monograph \textit{The Philosophy of Russell}, completed in 1965 and published in 1988~\cite{Jin1988}, and with his later English manuscript `On the Essence of Russell's Neutral Monism'~\cite{JinNeutralMonism}. The 1960 essay states the problem in a compressed form. The later writings give the more systematic target: Russell's attempt to replace inferred objective entities with logical constructions out of sense-data, and more generally his tendency to weaken the independent reality of the objects of knowledge. The Carroll discussion in 1960 is thus not an isolated exercise. It is the opening of Chin's later critique of Russell's treatment of inference, construction, and objectivity.

\section{Carroll's Paradox and Chin's Doctrine of Inference}

In the first two sections of `On ``Therefore''', Chin presents the doctrine of inference that underlies his response to Carroll. The doctrine has three components.

The first concerns implication. Chin identifies ``therefore'' with inference: inference is the transition from asserting a premise to asserting a conclusion. For this transition to be genuine, there must be a well-founded implication between premise and conclusion. This implication is objective. It holds or fails to hold independently of whether any particular reasoner recognises it or asserts it.

The second component concerns the truth of the premise. Inference requires not only that the premise imply the conclusion, but also that the premise be asserted as true. For Chin, however, the ground of this assertion is not merely the reasoner's acceptance. It is the objective necessity or factual basis that makes the premise true.

The third component concerns the occurrence of inference. Even when the premise is true and the implication holds, the inference does not happen automatically. It occurs only when a reasoner in fact makes the transition from asserting the premise to asserting the conclusion. Chin describes this as depending on ``historical conditions''~\cite[p.~49]{Liu2017}. This phrase is easy to misunderstand (for its layered meaning, see Section~\ref{sec:historical-background}). It should not be taken to mean that logical validity varies with historical circumstances. Rather, it marks the fact that inference is an act situated in a concrete epistemic context. History affects the occurrence and uptake of inference, not the validity of implication itself.

The point may be put as follows. Implication belongs to the objective order of logical and factual relations. Assertion belongs to the order of acts performed by thinkers. Inference stands at the junction of the two. It is not reducible to the act, because it is governed by objective implication; but it is not exhausted by the objective relation, because it must be performed by a finite reasoner under determinate conditions. This is the source of both the power and the tension of Chin's account.

\subsection{The Bolzano--Carroll Regress}

Lewis Carroll's 1895 paper presents a dialogue between Achilles and the Tortoise about a simple deductive argument. From (A) things equal to the same thing are equal to each other, and (B) the two sides of a particular triangle are equal to the same thing, one is supposed to conclude (Z) that those two sides are equal to each other.\footnote{Carroll's argument is presented entirely in natural language, as is Chin's discussion in `On ``Therefore'''. The notation $A \wedge B \rightarrow Z$ introduced later in this section is a modern reconstruction introduced solely for expository clarity; it should not be taken to represent the symbolic vocabulary of either Carroll or Chin. Likewise, Russell's $p \supset q$ follows the notation of \textit{The Principles of Mathematics}~\cite[\S\,38]{Russell1903} rather than any formalism employed by Chin.} The Tortoise accepts A and B but refuses to accept Z. He then demands that Achilles add a further conditional, C, saying that if A and B are true then Z must be true. Once Achilles adds C, the Tortoise demands another conditional, D, and the process continues indefinitely.

Scholars have interpreted the regress in several ways: as exposing a gap between logical consequence and psychological acceptance, as an attack on psychologism, or as a puzzle about rule-following~\cite{Marion2016, Philie2007, Imholtz2016}. Chin's reading is distinctive but not eccentric. He takes Carroll to be demonstrating that formal logic alone cannot compel someone to perform an inference if that person refuses to make the transition. In this sense, inference is ``logically difficult''~\cite[p.~34]{Liu2018}. This is what Chin means by calling Carroll's paper an ``attack on inference.''

Following Marion's interpretation~\cite[p.~53]{Marion2016}, a structurally similar regress can be identified in Bolzano's \textit{Theory of Science} (1837). In \S\,199, Bolzano argues that if a linking proposition expressing the inferential connection between grounds and consequent must itself be included in the full ground, then a further linking proposition will be required for that extended ground, and so on without limit. Engel, following Marion, has shown that Bolzano reached this conclusion while asking whether someone justified in accepting a premise is thereby justified in accepting the conclusion: treating the validity of inference as an explicit additional premise generates the same regress that Carroll later dramatised~\cite{Engel2016}.\footnote{Following Marion's recommendation for historical accuracy, we refer to this as the Bolzano--Carroll regress. For discussions of Bolzano's formulation and its relation to Carroll, see Engel~\cite{Engel2016}. Chen's~\cite{ChenM2026} recent claim that Wilson~\cite{Wilson1926} describes an argument bearing a strong structural resemblance to Carroll's. If confirmed, this would broaden the genealogy of the paradox beyond Bolzano and Carroll. As the claim remains new and requires further historical verification, we retain Marion's convention here without prejudice to Chen's findings.}

\subsection{Chin's Supplement to Carroll}

Chin's reconstruction of the regress is unusual because he thinks Carroll has not gone far enough. Carroll focuses on implication: even if A and B imply Z, the Tortoise can demand a further proposition licensing the transition. Chin observes that a similar regress can be generated on the side of the premises themselves. If the Tortoise refuses not only to move from A and B to Z but also to acknowledge that A and B are true, then Achilles may be tempted to add another proposition saying that A and B are true. Yet that new proposition would itself require assertion, and the same difficulty would arise again. Thus the regress threatens inference from two directions: from the side of implication and from the side of premise-assertion.

This supplement shows that Chin is not merely repeating the standard story of Carroll's paradox. He sees that inference requires both objective implication and asserted premises. If either element is turned into a further asserted premise, the inferential act is displaced and the regress begins again. The problem is not solved by adding more propositions. It is solved by distinguishing the role of propositions from the act of asserting and using them.

\subsection{Chin's Refutation: Proposition and Assertion}

Chin's response is to deny that the Tortoise's demand reveals a genuine defect in inference. The demand is misplaced because it treats the objective ground of inference as though it had to be added as another asserted premise. Once implication is confused with premise-assertion, every inference can be made to appear incomplete. But implication is not made true by being asserted. If $A \wedge B$ really implies $Z$, then the relation holds independently of whether Achilles or the Tortoise formulates it as another proposition. Likewise, if A and B are true, their truth does not depend on a further assertion that they are true.

This distinction between propositional content and assertion has an obvious place in the tradition that shaped Chin. It is closely related to Russell's distinction in \S\,38 of \textit{The Principles of Mathematics} between implication and inference.\footnote{A more distant formal ancestor is Frege's distinction between judgment and content in the \textit{Begriffsschrift}, expressed by the judgment stroke; see~\cite{vanderSchaar2013}.} Russell treats implication as an objective relation between unasserted propositions, while inference is an act involving the assertion of premises and conclusion~\cite[\S\,38]{Russell1903}. 

In this local respect, Chin's response to Carroll is Russellian. The Carroll discussion however supplies Chin with a point of resistance: if inference cannot be reduced to the assertion of further premises, then it must have an objective status that cannot be replaced by an endless sequence of psychological or linguistic acts. This is why the paradox becomes useful for Chin's later critique of Russell himself.

\section{Chin's Transformation of Assertion and the Two Faces of Inference}
\label{sec:assertion-two-faces}

The crucial passage in `On ``Therefore''' is not merely a parallel to Russell; it is an explicit engagement with Russell. Chin refers to \textit{The Principles of Mathematics}, \S\,38, and notes that Russell had already tried to strip assertion of its psychological component. Chin accepts the importance of this move, for it allows one to distinguish what is asserted from the act of asserting it, and hence to block the Carrollian regress. Yet Chin immediately marks his distance from Russell: his own account does not strip assertion of its cognitive component. Assertion is not merely the addition of a sign before a proposition, nor merely a private psychological episode. It is a cognitive act in which a thinker takes a proposition to be true in relation to an objective situation.

This point also clarifies Chin's use of the term ``objectivism.'' Chin's criticism of Russell's objectivism is not a criticism of objectivity as such. It is a criticism of a specific abstraction: the treatment of inference as the act of an isolated individual knower, detached from historical conditions and from the objective situation of cognition.\footnote{``Objectivism'' translates Chin's \emph{keguan zhuyi}. In the context of `On ``Therefore''', the term is polemical. It denotes an allegedly neutral and individualistic standpoint, not the realist objectivity of inference that Chin seeks to preserve.} For Chin, this abstraction weakens rather than secures the realist credentials of Russell's account, because it separates the act of assertion from the field of facts to which cognition is responsible.

The cognitive component does not give Chin a further logical advantage over Russell in answering Carroll. At the strictly formal level, the Russellian distinction between proposition and assertion is already sufficient to block the regress: implication need not be converted into another asserted premise. Chin's additional point does a different kind of work. It becomes important when the problem is no longer the formal account of inference, but the epistemic grounding of inference. Russell's distinction can prevent the Tortoise from turning implication into an infinite series of premises; it cannot, by itself, explain how inference remains connected with objective facts once Russell's later constructional programme begins to replace inferred entities by logical constructions. This is where Chin's insistence on the cognitive component of assertion becomes philosophically decisive.

This also helps clarify the relation between deductive and non-deductive inference. Chin should not be credited with a unified formal theory of both. The texts do not support such an attribution. Deductive and non-deductive inferences differ in structure and in standards of justification: the former are answerable to relations of implication, while the latter are answerable to facts, evidence, and historically situated scientific practice. What Chin offers is not a single technical account of inference in general, but a broader realist constraint on inference. A genuine cognitive transition must not be enclosed within the isolated subject's acts of assertion, construction, or psychological transition.

This point is especially important because Russell's own use of ``inference'' shifts across contexts. In \textit{The Principles of Mathematics}, the issue is assertion, implication, and the transition from premises to conclusion. In the constructional programme, the issue is largely non-deductive or epistemological inference from sense-data to physical objects or other objective entities. Chin was well aware of this difference. In \textit{The Philosophy of Russell}, he first considers the issue within propositional deductive systems and then turns to the broader problem of inference and construction. His criticism of Russell is therefore two-sided. He accepts the Russellian distinction between proposition and assertion against Carroll, but he resists the later Russellian tendency to replace inference to objective entities by logical construction. The common thread is not a formal theory of inference in general, but a realist insistence that inference must remain answerable to objective grounds.

\section{From Logical Construction to the Critique of Russell}
\label{sec:construction-russell}

The critique of construction is the point at which Chin's discussion of Carroll becomes more than a local response to a logical puzzle. Carroll's Tortoise threatens deductive inference by turning every rule or implication into a further premise. Russell's constructional programme, at least in the cases most relevant to Chin's criticism, threatens non-deductive inference by making inference to objective entities appear unnecessary. Chin does not collapse these two problems into one. What he sees, however, is that both can displace inference from its objective role. Carroll makes inference look impossible unless it is supplemented by an endless series of assertions; Russell makes inference look dispensable once construction has replaced the inferred object.

Chin's resistance to both moves is realist in spirit. Inference may take different forms in different contexts, but it must not be reduced either to the assertion of further propositions or to the psychological construction of subjective materials. This is why Chin's appeal to ``historical facts'' in `On ``Therefore''' should not be read as a lapse into relativism or mere ideological vocabulary. It expresses a realist conviction. Inference takes place under historical conditions because reasoning is performed by agents situated in a world of objective facts. Those historical conditions do not create logical validity. They are the field in which objective implication becomes operative in deductive reasoning, and in which factual constraint and scientific warrant become operative in non-deductive reasoning.

In \textit{The Philosophy of Russell}, Chin argues that Russell's constructional programme remains too close to the British empiricist tradition of Berkeley and Hume~\cite{Jin1988}. By beginning with sense-data understood as subjective materials, Russell weakens the link between cognition and the objective structure of the world. Sense-data are subjective and episodic. Logical constructions built out of them may be formally elegant, but they do not by themselves yield the objective facts that realism requires. Chen Bo has argued that Russell's logical atomism presupposes an isomorphism between language and world that Russell never fully grounds~\cite[pp.~944--945]{ChenB2019b}. A similar worry has been expressed in Russell scholarship: Russell's constructional method can clarify structural relations among the data, but it remains controversial whether such structures secure the independent reality of the physical world rather than merely redescribe patterns within experience~\cite{DemopoulosFriedman1985,LinskyLogicalConstruction}. This is where Chin's disagreement with Russell becomes explicit. If the relation between cognition and objective entities is mediated only by constructions out of subjective data, then the objectivity of inference is weakened. The bridge to the world has been replaced by an internal arrangement of experiences.

In this respect, Chen Bo is right to stress that Chin's account is not naive realism~\cite[pp.~937--944]{ChenB2019b}. Facts are not simply external items waiting to be mirrored. They involve both ``hardness'' and ``softness'': objectivity on the side of what is given, and cognitive organisation on the side of the knowing subject. This makes Chin's position more complex than Russell's early atomism and more realist than Russell's later sense-data construction. What Chin wants to preserve is the hard objective anchorage of inference without denying that inference is performed by finite cognitive subjects.

\section{Neutral Monism and the Reality of Objects}

Chin's later English manuscript `On the Essence of Russell's Neutral Monism' extends the same criticism from inference to the objects of knowledge. The paper was written after `On ``Therefore''' and around the same period as \textit{The Philosophy of Russell}. It is therefore important evidence that Chin's criticism of Russell in the 1960s was not an isolated polemical episode but a sustained philosophical project.

The central claim of that manuscript is that Russell's neutral monism is, in essence, a sense-data theory. Chin argues that Russell's attempt to construct both mind and matter out of neutral materials continues the Humean and Machian tendency to dissolve objects into collections of impressions, sensations, or events. What is lost in this dissolution is precisely what realism requires: the independent existence and relative stability of things. An apple is not merely a bundle of colour, shape, smell, and tactile impressions; a table is not merely a logical arrangement of events; a physical object is not exhausted by the sensory materials through which it is known. Chin calls this tendency the disembodiment of things.

The point is not that Chin rejects scientific description. On the contrary, he repeatedly distinguishes levels or spheres of discourse. Microphysical description may legitimately describe bodies in terms of particles, fields, or waves; ordinary perceptual and experimental practice legitimately treats tables, stones, feet, instruments, liquids, and gases as macroscopic objects. The mistake is to use one level to cancel another. If the microscopic description of a table is used to deny the reality of the ordinary table, or if the construction of physical objects out of sense-data is used to deny the independent existence of those objects, then scientific abstraction has been turned into metaphysical reduction.

Chin further appeals to Einstein in order to reject the idea that modern physics itself supports Russell's neutral monism. Einstein's realist orientation toward physical theory supports Chin's central claim that scientific cognition depends on a mind-independent world, even though such cognition is always mediated by historically situated theories, instruments, and practices. Physical theory may transform our account of that world, but it cannot abolish the independent reality of the world without undermining its own experimental basis. Experiments occur through instruments, bodies, materials, and observable events. If all these are reduced to subjective data or constructions, then the scientific theory that depends on them loses its foothold. Chin's criticism of neutral monism therefore parallels his criticism of construction: a method designed to clarify knowledge ends by weakening the reality of what is known.

This later argument helps us understand the philosophical core of `On ``Therefore'''. In both contexts Chin is concerned with a bridge. In `On ``Therefore''', the bridge is inference: the movement from asserted premises to asserted conclusions under objective logical and factual constraints. In the critique of neutral monism, the bridge is scientific cognition: the movement from experience and theory to a world of objects that exists independently of experience. 

This parallel is not accidental. It reveals a single structural concern in Chin's philosophy. Whether the issue is inference in `On ``Therefore''' or scientific cognition in the critique of neutral monism, Chin is trying to preserve a relation between finite cognitive activity and an objective reality that exceeds it. The knowing subject is historically situated and never operates from a view from nowhere; yet cognition must not be reduced to the subject's own materials, constructions, or psychological transitions. In this sense, the critique of neutral monism extends the same realist impulse that governs Chin's treatment of inference.

\section{Historical Facts, the Possibility of Thinking, and Sher's Foundational Holism}

The constructive upshot of Chin's discussion is stated in a key passage from `On ``Therefore''':

\begin{quote}
This transition [``therefore''] is something remarkable. It is a transition from asserting the premise to asserting the conclusion, but at the same time it is a bridge between ``logic'' and history. It crosses both the field of historical facts and the field of the possibility of thinking. Although ``therefore'' accepts those possible conditions of thinking, it is not entirely dominated by these conditions. It is determined by historical conditions, and it reflects those possible laws of thinking under historical conditions.~\cite[pp.~47--48]{Liu2017}
\end{quote}

In this passage, Chin identifies two fields. The first is the field of historical facts: the objective structure of reality in which reasoning is situated. The second is the field of the possibility of thinking: the formal and normative conditions under which inference is possible. Inference crosses both. It is an act of thought, but not a merely mental event. It is historically situated, but not historically relative in its validity.

This dual structure also connects Chin's discussion of Carroll with the logocentric predicament. In the `Prolegomena' papers of 1927, Chin had already asked how logic can have a starting point. Since philosophy depends on argument, and argument presupposes logic, logic appears to be both necessary and impossible to ground without circularity~\cite[pp.~37--38]{Chin1927a}. He described the starting points of philosophy as resting on ``prejudice,'' in the sense of unavoidable pre-theoretical commitment~\cite[p.~28]{Chin1927a}. His early answer, that it is ``convenient to believe in logic''~\cite[p.~23]{Chin1927b}, was deliberately unsatisfactory because it used a notion whose analysis already presupposed logical norms.

In 1931, Chin studied with Sheffer during his visit to Harvard University in Cambridge, Massachusetts. A decade later, in the epistemological treatise \textit{Theory of Knowledge}, he returned explicitly to Sheffer's ``logocentric predicament''~\cite[vol.~5, pp.~21--22]{Jin2014}. In \textit{Theory of Knowledge}, he formulates the point in his own terms: ``the idea of logic is not derived from logical systems. Suppose that one has not at all the idea of logic; then even if one pores over logical systems one would not get the idea of logic. This is the so-called `logic assumes logic' ''~\cite[p.~574]{Jin1983}. The parallel with Wittgenstein's remark that ``logic must take care of itself'' is evident~\cite[props.~5.473--5.4731]{Wittgenstein1961}.

Carroll's paradox is one manifestation of this predicament. If inference requires a further logical licence in order to begin, and if each licence requires another, logic cannot ground itself from within. Chin's answer is not to give a non-circular foundation in the strict sense. Rather, he shifts the ground. Inference is grounded neither in formal logic alone nor in historical fact alone, but in the crossing of the two. This is why the appeal to history is not an abandonment of realism. It is Chin's way of insisting that logic is not a self-enclosed calculus detached from the objective world.

This position has a cautious structural relationship with Gila Sher's foundational holism. Sher challenges the traditional opposition between foundationalism and coherentism, arguing that knowledge, including logical knowledge, can be grounded in reality without that grounding following a single linear hierarchy~\cite[pp.~353--355]{Sher2011}. For logic in particular, Sher writes:

\begin{quote}
As a theory of the formal structure of ``real'' objects (configurations of objects), logic is anchored in reality; as a theory of the formal structure of our thought of objects, logic is anchored in the mind.~\cite[p.~272]{Sher2016}
\end{quote}

In Sher's account, the worldly grounding of logic is technically secured by the criterion that logical constants are operations maximally invariant under permutations of the domain of objects~\cite[pp.~160--167]{Sher2013}. This invariance criterion, developed from Tarski's 1966 lecture ``What Are Logical Notions?''~\cite{Tarski1986}, gives Sher's account formal precision. Logical constants denote formal features of the world in a mathematically articulated sense. The philosophical force of this criterion is that logicality is not a mere product of linguistic convention. A term counts as logical because it is insensitive to the identities of particular objects and tracks only the structural relations that remain invariant across transformations of the domain.

Chin did not possess this apparatus and should not be described as anticipating Sher's theory. The comparison is heuristic rather than genealogical. The parallel lies at a more general level. Both Chin and Sher reject the idea that logic can be grounded in the mind alone or in the world alone. For Sher, logic is formal but world-involving. For Chin, inference is an act of thought but not a merely subjective or psychological occurrence. In both cases, the objectivity of logic requires a relation between the structure of thought and the structure of reality.

The differences are equally important. Sher's theory is formally developed; Chin's ``field of historical facts'' is a programmatic metaphysical category. The contrast between them should not be treated as a simple ranking of strengths and weaknesses. Sher's formal thinness is the price of model-theoretic precision: the invariance criterion identifies logicality by abstracting from the identities of particular objects. Chin's ontological thickness raises a different question, one that the invariance criterion by itself does not answer: how are the formal structures tracked by logic related to historically situated acts of knowing and inferring? Conversely, Chin lacks the mathematical precision that gives Sher's account its contemporary power. The two views are therefore not competing answers to exactly the same question, but two poles of a problem that neither alone resolves. What matters for the present paper is that Chin's appeal to historical facts is not a concession to relativism. It is the realist side of his account of inference.


\section{Conclusion: Two Ideals, or One Ideal in Two Registers?}\label{sec:conclusion}

Hao Wang famously suggested that Chin pursued two fundamentally different philosophical ideals over the course of his career~\cite{Wang1984}. Before 1949, Wang argues, Chin's ideal was the construction of a comprehensive neorealist philosophical system, grounded in logic and shaped by the Russell--Moore tradition. This project was expressed in \textit{Logic} (1936), \textit{On Dao} (1940), and the draft of \textit{Theory of Knowledge} completed in 1948. After 1949, Chin's public philosophical language changed, and his writings appeared under the vocabulary of dialectical materialism.

Wang's interpretation is not wrong at the level of form and vocabulary. However, the present paper challenges Wang's thesis at a deeper level of philosophical content. The evidence now runs through several stages. In the 1927 `Prolegomena' papers, Chin asked how logic can have a starting point without presupposing logic. In `On ``Therefore''', he returns to Russell's \textit{The Principles of Mathematics} and to the problem of assertion in order to defend the objectivity of inference. In his 1962 response to Zhou Liquan, he distinguishes more carefully between implication, inference, inference form, and concrete inferential activity. In \textit{The Philosophy of Russell}, he criticises Russell's attempt to replace inference with construction. In the English manuscript on neutral monism, he extends the same criticism to Russell's sense-data theory and to the reality of objects. The vocabulary changes; the organising concern remains.

The concern that runs through these texts is the objectivity of logic, inference, and knowledge. How can inference be an act of thought without becoming merely psychological? How can logic be objective without becoming a self-enclosed formal calculus detached from historical facts? How can scientific cognition be historically situated without losing its claim to objects that exist independently of cognition? These questions persist across the apparent divide in Chin's career. The difference between Wang's reading and ours is therefore not a disagreement about the external facts of Chin's career, but a difference in philosophical emphasis: Wang sees rupture in system and vocabulary; we see continuity in the problem that continued to organise Chin's thinking.

The Carroll section of `On ``Therefore''' is therefore not a marginal episode. It is the point at which Chin's long-standing concern with the foundation of logic becomes tied to his later critique of Russell. Locally, Chin's answer to Carroll is Russellian: proposition must be distinguished from assertion, implication from inferential performance. Philosophically, however, Chin uses this answer against what he takes to be Russell's later weakening of inference. If inference can be replaced by constructions out of sense-data understood as subjective materials, then the bridge between thought and objective fact is broken. Chin's insistence on historical facts is his way of restoring that bridge.

Chin did not solve Carroll's paradox once and for all, nor did he provide a modern inferentialist account of logical normativity. His contribution is different. He recognised that the paradox matters because it concerns the status of inference itself. Inference is not merely a relation between propositions in the abstract, nor merely a psychological transition among beliefs. It is the place where the mind's formal capacity for thought meets a factual world that exists independently of that thought. This is why Chin's engagement with Carroll and Russell deserves a place not only in the history of modern Chinese philosophy but also in the broader history of reflection on the objectivity of logic.

\subsection*{Acknowledgements}

Thanks go to Amirouche Moktefi, Ahti-Veikko Pietarinen, and Liying Zhang for their encouragement and advice in preparing this paper. This work is supported by the National Social Science Fund of China (\textit{A Study of Chin Yueh-Lin's Academic Thought from the Perspective of Logic and Philosophical System Construction}, 24AZX023).

\bibliographystyle{abbrv}
\bibliography{jin_logic}

\end{document}
